\section{The height equation}
\label{sec:scalar-core}

The geometric reduction leaves two things to choose: the rotational
metric, that is the slope $w$, and the nonradial height $z$. This
section derives the equation that relates them and extracts from it
the information that guides both choices.
\begin{itemize}
\item Section~\ref{subsec:residual} writes the curvature of the
 remainder metric as the exterior derivative of a connection form.
 This gives the equation, its linearization in divergence form, and a
 conservation law.
\item Section~\ref{subsec:monotone} shows that a flat end at which $w$
 is monotone carries no infinitesimal bending. This is why the profile
 must oscillate.
\item Section~\ref{sec:worked-example} computes, in a finite model, the
 two different resonance sets mentioned in the introduction.
\item Section~\ref{subsec:normalization} normalizes the equation in
 the small parameter, and Section~\ref{subsec:first-correction}
 introduces the first of the two residual corrections.
\end{itemize}

\subsection{The curvature residual and its derivative}
\label{subsec:residual}

Let $h$ be a radial function with $h'=w$, so that $g$ in
\eqref{eq:intro-metric} is the metric of its graph and $E=1+w^2$. For
a height $z(r,\theta)$ we use the abbreviations
\begin{equation}\label{eq:height-jet-notation}
 k=z_r,\qquad p=\frac{z_\theta}{r},\qquad
 D=E(1-p^2)-k^2,\qquad Q=k+Ep_\theta .
\end{equation}
In the coordinates $(r,\theta)$ the remainder metric
$\bar g=g-dz\otimes dz$ has determinant $r^2D$, and
$1-|dz|_g^2=D/E$. Hence
\[
 \bar g\ \text{is positive definite}
 \iff D>0\iff |dz|_g<1 .
\]
We call $D$ the \emph{discriminant}. When $D>0$, an oriented
$\bar g$-orthonormal coframe is
\[
 \eta_1=\frac{\sqrt D}{\sqrt{1-p^2}}\,dr,\qquad
 \eta_2=r\sqrt{1-p^2}\,d\theta-\frac{kp}{\sqrt{1-p^2}}\,dr .
\]
For the radial height $z=h$ it reduces to the polar coframe
$dr,\ r\,d\theta$ of the plane, whose connection form is $d\theta$.

All adjoints of scalar operators in this paper are taken with respect
to $dr\,d\theta$, and source terms are densities for this measure. The
measure will be written out whenever coordinates are changed.

\begin{proposition}[Connection form, residual, linearization]
\label{prop:scalar-core}
Assume $D>0$.
\begin{enumerate}
\item The connection form $\alpha$ of the coframe $(\eta_1,\eta_2)$,
 defined by $d\eta_1=\alpha\wedge\eta_2$ and
 $d\eta_2=-\alpha\wedge\eta_1$, is
 \begin{equation}\label{eq:conn-one-form}
  \alpha=D^{-1/2}\left(d\theta+\frac{k}{1-p^2}\,dp\right).
 \end{equation}
\item Define the \emph{curvature residual}
 $\mathcal N_h(z)=-\sqrt D\,K_{\bar g}$. Then
 $d\alpha=r\,\mathcal N_h(z)\,dr\wedge d\theta$ and
 \begin{equation}\label{eq:cart-curvature-scalar}
  r\,\mathcal N_h(z)=D^{-3/2}
  \left[Qk_r-Er\,p_r^2-\frac{E'}2\,(1-p^2+kp_\theta)\right].
 \end{equation}
 In particular $\bar g$ is flat if and only if
 \begin{equation}\label{eq:darboux-numerator}
  Qk_r-Er\,p_r^2-\frac{E'}2\,(1-p^2+kp_\theta)=0 .
 \end{equation}
\item For fixed $w$, the derivative of the residual in the direction
 $\eta$ is in divergence form,
 \begin{equation}\label{eq:cart-true-derivative}
  r\,D_z\mathcal N_h(z)\,\eta
   =\partial_r(a\eta_r+b\eta_\theta)
       +\partial_\theta(b\eta_r+c\eta_\theta),
 \end{equation}
 with coefficients
 \begin{equation}\label{eq:linear-coefficients}
  a=\frac{Q}{D^{3/2}},\qquad b=-\frac{Ep_r}{D^{3/2}},\qquad
  c=\frac{Ek_r-E'k/2}{r\,D^{3/2}} .
 \end{equation}
\item The determinant of the coefficient matrix is
 \begin{equation}\label{eq:cart-determinant-identity}
  ac-b^2=\frac{K_0}{D^2}+\frac{E}{D^{3/2}}\,\mathcal N_h(z),
  \qquad K_0=\frac{ww'}r .
 \end{equation}
\end{enumerate}
\end{proposition}

\begin{proof}
Equality of mixed derivatives of $z$ gives the compatibility relation
\begin{equation}\label{eq:height-compatibility}
 r\,p_r=k_\theta-p .
\end{equation}

\step{1: the connection form}
Write $\alpha=\alpha_r\,dr+\alpha_\theta\,d\theta$. The second
structure equation determines $\alpha_\theta$:
\[
 \alpha_\theta\,\frac{\sqrt D}{\sqrt{1-p^2}}
 =\partial_r\bigl(r\sqrt{1-p^2}\bigr)
  +\partial_\theta\!\left(\frac{kp}{\sqrt{1-p^2}}\right)
 =\frac1{\sqrt{1-p^2}}+\frac{kp_\theta}{(1-p^2)^{3/2}},
\]
where \eqref{eq:height-compatibility} was used. This is the
$d\theta$-component of \eqref{eq:conn-one-form}. The first structure
equation then reads
$\alpha_r\,r\sqrt{1-p^2}+\alpha_\theta\,kp/\sqrt{1-p^2}
 =-\partial_\theta(\sqrt D/\sqrt{1-p^2})$,
and with $D_\theta=-2Epp_\theta-2kk_\theta$ and
\eqref{eq:height-compatibility} it is satisfied by
$\alpha_r=D^{-1/2}kp_r/(1-p^2)$.

\step{2: the residual}
For an orthonormal coframe $d\alpha=-K_{\bar g}\,\eta_1\wedge\eta_2$,
and here $\eta_1\wedge\eta_2=r\sqrt D\,dr\wedge d\theta$; this is the
relation between $d\alpha$ and $\mathcal N_h$. To compute $d\alpha$,
note that
\[
 \alpha=D^{-1/2}\,d\theta+B_1\,dp,\qquad
 B_1=\frac{k}{(1-p^2)\sqrt D},
\]
and that $D_r=E'(1-p^2)-2Epp_r-2kk_r$. The coefficient of
$dr\wedge d\theta$ in $D^{3/2}\,d\alpha$ is
\[
 -\tfrac12D_r+\frac{D}{1-p^2}(k_rp_\theta-k_\theta p_r)
 -\frac{k}{2(1-p^2)}(D_rp_\theta-D_\theta p_r).
\]
Since $D_rp_\theta-D_\theta p_r
 =E'(1-p^2)p_\theta-2k(k_rp_\theta-k_\theta p_r)$ and $D+k^2=E(1-p^2)$,
this equals
\[
 -\frac{E'}2(1-p^2+kp_\theta)+k_r(k+Ep_\theta)-Ep_r(k_\theta-p),
\]
which is the bracket in \eqref{eq:cart-curvature-scalar} by
\eqref{eq:height-compatibility}. As a check, the Christoffel symbols
$\Gamma^r_{rr}=E'/2E$, $\Gamma^r_{\theta\theta}=-r/E$,
$\Gamma^\theta_{r\theta}=1/r$ and the curvature
\begin{equation}\label{eq:radial-curvature}
 K_g=\frac{E'}{2rE^2}
\end{equation}
of $g$ give the same expression from the invariant form
\eqref{eq:intro-darboux} of the Darboux equation.

\step{3: the variation}
Vary $z$ by $\eta$, so that $\delta k=\eta_r$ and
$\delta p=\eta_\theta/r$. Since
$\delta\alpha=\delta(D^{-1/2})\,d\theta+\delta B_1\,dp-\delta p\,dB_1
 +d(B_1\,\delta p)$,
and since the partial derivatives of $B_1$ are
$\partial_kB_1=ED^{-3/2}$ and $\partial_EB_1=-k/(2D^{3/2})$ while its
$p$-derivative cancels, a short computation gives
\begin{equation}\label{eq:conn-exact-variation}
 \delta\alpha=d(B_1\,\delta p)
  +(a\eta_r+b\eta_\theta)\,d\theta
  -(b\eta_r+c\eta_\theta)\,dr .
\end{equation}
Taking the exterior derivative removes the exact term and proves
\eqref{eq:cart-true-derivative}. The formula also identifies
$a\eta_r+b\eta_\theta$ as the flux of the linearized equation through
a circle $r=\text{const}$; we call it the \emph{normal flux}.

\step{4: the determinant}
Let $N$ be the bracket in \eqref{eq:cart-curvature-scalar}. From
$kQ+D=E(1-p^2+kp_\theta)$ one finds
\[
 ac-b^2=\frac E{rD^3}\left(N+\frac{E'D}{2E}\right),
\]
and $N=rD^{3/2}\mathcal N_h(z)$, $E'=2ww'$ give
\eqref{eq:cart-determinant-identity}.
\end{proof}

\paragraph{The linearization at the radial height.}
At $z=h$ we have $k=w$, $p=0$, $D=1$, and the coefficients are $a=w$,
$b=0$, $c=w'/r$. The linearized operator is therefore
\begin{equation}\label{eq:intro-linearization}
 \mathscr L_wu=(wu_r)_r+\frac{w'}r\,u_{\theta\theta}.
\end{equation}
For $w>0$ it is elliptic where $w'>0$ and hyperbolic where $w'<0$; by
\eqref{eq:radial-curvature} these are the regions of positive and
negative Gaussian curvature of $g$. A simple zero of $w'$ is called a
\emph{turning circle}. Note that it is the derivative of the slope
that changes sign there, not the slope itself: $w$ stays positive.

\subsection{Why the profile must oscillate at a flat end}
\label{subsec:monotone}

It is tempting to take for $w$ a simple bump, to solve
$\mathscr L_wu=0$ in the middle of the annulus, and to flatten the
ends afterwards. The following observation shows that a nonzero
infinitesimal height cannot survive this. Separating variables,
$u=f(r)\cos(j\theta)$, the equation $\mathscr L_wu=0$ becomes
$(wf')'-j^2w'f/r=0$.

\begin{proposition}[The monotone-end obstruction]
\label{prop:obs-monotone-tail}
Let $w>0$ be $C^1$ on $(r_-,r_+)$ with $w(r_-)=0$ and $w'\ge0$ near
$r_-$. Let $j\ge1$ and let $f$ be a solution of
\[
 (wf')'-j^2\,\frac{w'}{r}\,f=0
\]
which extends continuously to $r_-$ with
\[
 f(r_-)=0,\qquad \lim_{r\downarrow r_-}w(r)f'(r)=0 .
\]
Then $f=0$. The same holds at the right endpoint if $w(r_+)=0$ and
$w'\le0$ near $r_+$.
\end{proposition}

\begin{proof}
Integrate the equation from $r_-$ and put
$f^*(r)=\max_{r_-\le t\le r}|f(t)|$. Since $w'\ge0$ and
$\int_{r_-}^rw'=w(r)$,
\[
 |f'(r)|=\frac{j^2}{w(r)}
       \left|\int_{r_-}^r\frac{w'(t)}t\,f(t)\,dt\right|
 \le\frac{j^2}{r_-}\,f^*(r).
\]
Integrating once more gives
$f^*(r)\le(j^2/r_-)\int_{r_-}^rf^*$, so $f^*=0$ near $r_-$ by
Gronwall's inequality, and $f=0$ everywhere by uniqueness for the
ordinary differential equation in the interior. At the right endpoint
integrate backwards and use $\int_r^{r_+}|w'|=w(r)$. The argument
applies to complex-valued $f$ as well.
\end{proof}

A nonconstant angular mode of a height with constant extensions
satisfies both boundary conditions of the proposition. Consequently a
positive profile with only finitely many changes of monotonicity cannot
carry a nonzero angular mode which is flat at the endpoints. Within
our construction the infinitely many turning circles are therefore
forced. Their spacings and the amplitude of $w$ between them will
decrease so fast that $w$ nevertheless extends smoothly by zero
(Figure~\ref{fig:profile-type}).

\begin{figure}[htbp]
\centering
\begin{tikzpicture}[x=1cm,y=1cm,>=Stealth,font=\small]
\draw[->] (0,0)--(11,0) node[right] {$r$};
\draw[thick,blue!70!black] (0,0)--(.4,0)
 .. controls (.6,0) and (.65,.15) .. (.8,.1)
 .. controls (.95,.05) and (1,.4) .. (1.2,.3)
 .. controls (1.4,.2) and (1.5,.85) .. (1.8,.65)
 .. controls (2.1,.45) and (2.2,1.4) .. (2.8,1.2)
 .. controls (3.5,1) and (3.5,2) .. (4.5,1.75)
 .. controls (5.5,1.4) and (6,2.1) .. (6.8,1.5)
 .. controls (7.5,.9) and (7.5,1.4) .. (8.0,.85)
 .. controls (8.5,.35) and (8.5,.8) .. (8.8,.4)
 .. controls (9.0,.16) and (9.1,.36) .. (9.3,.16)
 .. controls (9.5,0) and (9.6,.12) .. (9.8,0)--(10.5,0);
\node[above] at (5.2,2.15) {positive profile $w$};
\node[below,align=center] at (1.4,-.2) {accumulating turns,\\flat extension};
\node[below,align=center] at (8.9,-.2) {accumulating turns,\\flat extension};
\draw[->,orange!70!black] (4.8,1.1)--(5.2,1.6);
\node[below] at (5.1,.9) {$w'<0$: hyperbolic};
\draw[->,teal!70!black] (3.1,.45)--(3.65,1.5);
\node[below] at (3.1,.4) {$w'>0$: elliptic};
\end{tikzpicture}
\caption{The shape of an admissible profile. Only a few representative
turns are drawn, and not to scale. Where $w$ rises the linearized
equation is elliptic, where it falls it is hyperbolic. Every derivative
of $w$ tends to zero at the two accumulation points.}
\label{fig:profile-type}
\end{figure}

\subsection{Two resonance conditions in a finite model}
\label{sec:worked-example}

Before constructing the profile we examine a model with exactly two
turning circles, in which everything can be computed. It is a guide to
the construction of Part~\ref{part:reference} and is not used in any
proof. Its purpose is to show that an inverse of $\mathscr L_w$
involves two different sets of small denominators.

\paragraph{The model.}
Put $\epsilon=\sqrt2/100$ and $\rho(X)=\epsilon^2X(1-X)$, and let
$r(X)$ solve, on $-1\le X\le2$,
\begin{equation}\label{eq:example-radius}
 r_{XX}=-\rho\,r,\qquad r(0)=1,\qquad r_X(0)=\frac1{10}.
\end{equation}
On this interval $|\rho|\le2\epsilon^2$, so the integral form of
\eqref{eq:example-radius} bounds $R=\max|r|$ by
$R\le1.2+4\epsilon^2R$. Hence $R<1.21$, $r>0.89$ and
$|r_X-\tfrac1{10}|\le4\epsilon^2R<0.001$. In particular $r$ is
increasing, and we may define a positive profile by
\[
 w(r(X))=r_X(X).
\]
Then $w'=dw/dr=-\rho\,r/r_X$, so $w$ has exactly two turning circles,
at $X=0$ and $X=1$, and the equation is hyperbolic between them. For an
angular mode $u=U(X)\cos(j\theta)$ the operator
\eqref{eq:intro-linearization} becomes, with $J=j\epsilon$,
\begin{equation}\label{eq:example-mode}
 L_jU=\frac1{r_X}\bigl(U_{XX}+J^2X(1-X)\,U\bigr).
\end{equation}

\paragraph{First denominator: characteristics.}
On $0<X<1$ the characteristics of $\mathscr L_w$ have slopes
$d\theta/dX=\pm\epsilon\sqrt{X(1-X)}$. Between the two turning circles
each of them turns by the \emph{travel angle}
\begin{equation}\label{eq:example-action}
 T=\epsilon\int_0^1\sqrt{X(1-X)}\,dX=\frac{\epsilon\pi}8 .
\end{equation}
The two families of characteristics translate a function on the circle
by $+T$ and by $-T$, and reconciling them requires inverting the
difference of the two translations,
\begin{equation}\label{eq:example-circle-operator}
 \mathcal D_Tf(\theta)=f(\theta+T)-f(\theta-T),\qquad
 \mathcal D_T\sin(j\theta)=2\sin(jT)\cos(j\theta).
\end{equation}
So the characteristic problem costs the divisor $\sin(jT)$. Here
$jT/\pi=j\sqrt2/800$ is a quadratic irrational multiple of $j$. If $m$
is the integer nearest to $j\sqrt2/800$, the integer
$2j^2-640000\,m^2$ is nonzero; factoring it, and using
$j\sqrt2/800+|m|\le j$ and $|\sin(\pi x)|\ge2\dist(x,\Z)$, we obtain
\begin{equation}\label{eq:example-diophantine-bound}
 |\sin(jT)|\ge\frac1{320000\,j}\qquad(j\ge1).
\end{equation}
The inverse of $\mathcal D_T$ exists on functions of mean zero and
loses one angular derivative. No choice of $T$ avoids a loss
altogether. If $T/\pi$ is rational, $\sin(jT)=0$ for infinitely many
even $j$. If it is irrational, apply the pigeonhole principle to the
fractional parts of $2kT/\pi$, $0\le k\le N$: two of them differ by
at most $1/N$, which gives an even $j\le2N$ with
$|\sin(jT)|\le\pi/N$. So the divisors are never bounded below, and
a bound like \eqref{eq:example-diophantine-bound} is the best one
can ask for.

\paragraph{Second denominator: matching.}
Now consider the boundary value problem for \eqref{eq:example-mode}.
The substitution $x=\sqrt J\,(X-\tfrac12)$ turns the homogeneous
equation into the harmonic oscillator equation
\begin{equation}\label{eq:example-oscillator}
 (-\partial_x^2+x^2)\,U=\frac J4\,U .
\end{equation}
As boundary condition at each end of the interval we require that
$(U,U_X)$ be proportional to the data of the solution of
\eqref{eq:example-oscillator} on the whole line which decays in that
direction. These are legitimate boundary conditions for
\eqref{eq:example-mode}; the radius need not be extended beyond
$[-1,2]$. The homogeneous problem then has a nonzero solution exactly
when one solution decays in both directions, that is, when $J/4$ is an
eigenvalue of the oscillator. The eigenvalues are the odd integers
$2m+1$ \cite[Section~8.3]{Teschl2009}, so the resonances are
\begin{equation}\label{eq:example-spectral-phase}
 \frac J4=2m+1,
 \qquad\text{equivalently}\qquad
 jT=\bigl(m+\tfrac12\bigr)\pi .
\end{equation}
Away from them the inhomogeneous problem is solved by the resolvent of
the oscillator: multiply the source by $-r_X/J$, extend it by zero
outside $|x|\le3\sqrt J/2$, apply the resolvent on the whole line, and
restrict. Thus the boundary value problem has precisely the
denominators \eqref{eq:example-spectral-phase}; this is an exact
statement, not a comparison of formal phases.

\begin{figure}[htbp]
\centering
\begin{tikzpicture}[x=1.5cm,y=.8cm,>=Stealth,font=\small]
\node[anchor=east] at (-.2,1.5) {characteristic return};
\node[anchor=east] at (-.2,0) {matching of decaying solutions};
\draw[->] (0,1.5)--(4.4,1.5);
\draw[->] (0,0)--(4.4,0) node[right] {$jT$};
\foreach \x/\t in {0/0,1/\pi,2/2\pi,3/3\pi,4/4\pi}{
 \fill[blue!70!black] (\x,1.5) circle (2pt);
 \node[below] at (\x,-.15) {$\t$};}
\foreach \x in {.5,1.5,2.5,3.5}{
 \fill[orange!80!black] (\x,0) circle (2pt);}
\node[above] at (2,1.7) {$\sin(jT)=0$};
\node[above] at (2,.2) {$jT=(m+\tfrac12)\pi$};
\end{tikzpicture}
\caption{The two resonance sets of the finite model, as functions of
the phase $jT$. Avoiding one row does not avoid the other. In the
infinite construction both must be kept at a controlled distance.}
\label{fig:two-denominators}
\end{figure}

\paragraph{Remark on the half-integer shift.}
The shift by $\pi/2$ between the two rows of
Figure~\ref{fig:two-denominators} has a local explanation. Near $X=0$
put $y=J^{2/3}X$; then
\[
 U_{yy}+yU=J^{-2/3}y^2\,U .
\]
On a fixed $y$-interval the solutions are, up to an error
$O(J^{-2/3})$ in value and first derivative, combinations of the Airy
functions $\Ai(-y)$ and $\Bi(-y)$; this follows by variation of
constants and Gronwall's inequality. The asymptotics of the Airy
functions carry a phase shift $\pi/4$ at each turning point
\cite[9.7.5--9.7.12]{DLMFAiry}, and two turning points give $\pi/2$.
This is only a heuristic; \eqref{eq:example-spectral-phase} was proved
by the exact oscillator computation.

\paragraph{What the model teaches.}
In the actual construction the travel angle of every hyperbolic
annulus will be a quadratic irrational multiple of $\pi$, which gives
a bound of the type \eqref{eq:example-diophantine-bound}; and the
matching problem, now between infinitely many pieces, will be
controlled by spectral gaps of the individual pieces together with a
choice of their relative heights.

\subsection{Normalization and the conservation law}
\label{subsec:normalization}

The nonlinear unknown is $z=h+qu$ with a small parameter $q$. Since
$\mathcal N_h(h)=0$, the residual of $z$ is divisible by $q$, and we
define the \emph{normalized residual}
\begin{equation}\label{eq:cart-removable-equation}
 \mathcal F_q(w,u)=
 \begin{cases}r\,\mathcal N_h(h+qu)/q,&q\ne0,\\[2pt]
 \mathscr L_wu,&q=0.
 \end{cases}
\end{equation}
It is smooth in $q$ through $q=0$, because
\begin{equation}\label{eq:normalized-integral}
 \mathcal F_q(w,u)=\int_0^1 r\,D_z\mathcal N_h(h+tqu)\,u\,dt .
\end{equation}
It depends on $h$ only through $w$, and it does not change when a
constant is added to $u$.

\paragraph{Symmetry class.}
From now on all heights are even in $\theta$ and invariant under
$\theta\mapsto\theta+\pi$. This class is preserved by products and by
$\mathcal F_q$. Its Fourier modes are $\cos(j\theta)$ with $j\ge0$
even.

\begin{lemma}[Conservation law]\label{lem:scalar-integral}
Let $w$ be flat at both endpoints and let $u$ have constant
extensions at both endpoints, with $D>0$ for $z=h+qu$. Then
\begin{equation}\label{eq:zero-total-integral}
 \int_{r_-}^{r_+}\av{\mathcal F_q(w,u)}\,dr=0 .
\end{equation}
The same holds for every derivative of $\mathcal F_q(w,u)$ with
respect to $u$, $q$ or $w$ in directions preserving these conditions.
\end{lemma}

\begin{proof}
By Proposition~\ref{prop:scalar-core}(ii) and Stokes' theorem,
$\int\!\!\int r\mathcal N_h(z)\,dr\,d\theta$ is the difference of the
periods $\oint\alpha$ over the two boundary circles. At a boundary
circle $w=0$ and $dz=0$ to infinite order, so $D=1$ and
$\alpha=d\theta$ there: both periods equal $2\pi$. Differentiating the
identity gives the last statement.
\end{proof}

\paragraph{The connection period.}
The angular mean of the equation is most transparent at the level of
the connection form. For a height $z=h+U$ define
\begin{equation}\label{eq:cart-mean-path-flux}
 \mathfrak F_w(U)=\av{\alpha_\theta}-1
 =\left\langle\frac{1+kp_\theta/(1-p^2)}{\sqrt D}\right\rangle-1,
\end{equation}
the normalized period of $\alpha$ over the circle of radius $r$. Since
$\partial_r\alpha_\theta-\partial_\theta\alpha_r=r\mathcal N_h$, its
radial derivative is $r\av{\mathcal N_h(h+U)}$, and it vanishes at the
endpoints. So the mean of the equation $\mathcal N_h(h+U)=0$ is
equivalent to
\[
 \mathfrak F_w(U)=0\qquad\text{for all }r .
\]
Formula~\eqref{eq:conn-exact-variation} gives the derivative of
$\mathfrak F_w$ with respect to a radial variation $\bar U(r)$ of the
height, and its full derivative at $U=0$:
\begin{equation}\label{eq:mean-period-derivative}
 \partial_{\bar U'}\mathfrak F_w(U)=\av a,
 \qquad D_U\mathfrak F_w(0)[U]=w\av{U_r}.
\end{equation}
The second formula says that, to first order, the period sees only the
mean of the height. This cancellation is the reason why the mean of
the solution will decay twice as fast as its oscillating part
(Section~\ref{sec:final-proof}).

\subsection{The first residual correction}
\label{subsec:first-correction}

By \eqref{eq:cart-determinant-identity}, the type of the true
linearization \eqref{eq:cart-true-derivative} is governed by
$K_0/D^2+ED^{-3/2}\mathcal N_h(z)$. At a solution the second term
vanishes and the type changes exactly where $w'$ does. Away from a
solution the second term moves the turning circles, which is
inconvenient: the circles would depend on the unknown. We remove that
term by changing the tangential coefficient $c$ only.

\begin{lemma}[Fixing the type-change set]\label{lem:type-completion}
On the set of heights with $D\ge\tfrac12$ and $Q\ge\tfrac78w>0$ put
\[
 B=\frac{E}{rQ},\qquad \bar c=c-B\,r\,\mathcal N_h(z),\qquad
 \bar L=\partial_r(a\partial_r+b\partial_\theta)
             +\partial_\theta(b\partial_r+\bar c\,\partial_\theta).
\]
Then
\begin{equation}\label{eq:cart-completed-normalized-derivative}
 a\bar c-b^2=\frac{K_0}{D^2},\qquad
 D_u\mathcal F_q(w,u)\,\eta=\bar L\eta
       +q\,\partial_\theta\bigl(B\,\mathcal A\,\eta_\theta\bigr),
 \qquad \mathcal A=\mathcal F_q(w,u),
\end{equation}
where $\bar L$ is formed with the coefficients of $z=h+qu$. The
operator $\bar L$ is symmetric for $dr\,d\theta$, annihilates
constants, and has the same normal flux $a\eta_r+b\eta_\theta$ as the
true derivative.
\end{lemma}

\begin{proof}
The determinant changes by $-aB\,r\mathcal N_h=-ED^{-3/2}\mathcal N_h$,
which cancels the unwanted term of
\eqref{eq:cart-determinant-identity}. The operator changes by
$-\partial_\theta(B\,r\mathcal N_h\,\partial_\theta)$, and
$r\mathcal N_h(h+qu)=q\mathcal A$. Only the tangential coefficient was
modified, which gives the remaining assertions.
\end{proof}

Thus $\bar L$ is elliptic exactly where $w'>0$ and hyperbolic exactly
where $w'<0$, whatever the height. The price is the term
$q\,\partial_\theta(B\mathcal A\eta_\theta)$, which is the product of
the residual $\mathcal A$ and the increment $\eta$. An error of this
bilinear form is exactly what Hamilton's version of the implicit
function theorem tolerates (Section~\ref{sec:solving-the-equation}).
This is the first residual correction. The second, which fixes the
characteristic return maps, is made in
Section~\ref{sec:fixed-geometry}.

\paragraph{Characteristics.}
For later reference we record the characteristic directions of
$\bar L$ on a component of $\{w'<0\}$:
\begin{equation}\label{eq:completed-characteristics}
 \frac{d\theta}{dr}=d\pm v,\qquad
 d=\frac ba=\frac{E(p-k_\theta)}{rQ},\qquad
 v=\frac{\sqrt{b^2-a\bar c}}{a}
  =\frac{\sqrt{-ww'/r}\,\sqrt D}{Q}.
\end{equation}
We call $v$ the \emph{characteristic speed}. At the radial height
$d=0$ and $v=\sqrt{-w'/(rw)}$. At a turning circle $v$ is a smooth
positive factor times the square root of the distance into the
hyperbolic side. The coordinates of
Section~\ref{sec:fixed-geometry} are built to respect this
degeneracy; treating all intervals as uniformly elliptic or uniformly
hyperbolic would lose the geometry of the problem.
